\documentclass{article}
\usepackage[utf8]{inputenc}
\usepackage[T1]{fontenc}
\usepackage{geometry}
\usepackage{graphicx}
\usepackage{float}
\usepackage{hyperref}
\usepackage[english,indonesian]{babel}

\title{Optimasi Pemrosesan Template LaTeX Menggunakan Pendekatan Berbasis AI}
\author{Academic Support Agent}
\date{\today}

\begin{document}
\maketitle

\begin{abstract}

\end{abstract}

\chapter{PENDAHULUAN}


\chapter{TINJAUAN PUSTAKA}


% Original content (commented out):
% %%%%%%%%%%%%%%%%%%%%%%%%%%
%%% CHARABIA TECHNIQUE %%%
%%%%%%%%%%%%%%%%%%%%%%%%%%
\documentclass{article}
\usepackage[utf8]{inputenc}
\usepackage[chorded]{songs}
\usepackage[top=0.5cm,bottom=0.5cm,left=5.5cm,right=5.5cm]{geometry}
\usepackage[dvipsnames,table,xcdraw]{xcolor}
\input{bis}
\begin{document}
\songcolumns{1}
\renewcommand\printchord[1]{{\small\sffamily\it\color{black!60!white}#1}}
\renewcommand{\lyricfont}{\sffamily\large}
\renewcommand{\notebgcolor}{black!13!white}
\renewcommand{\snumbgcolor}{black!20!white}
%\selectlanguage{french}
\begin{songs}{}



%%%%%%%%%%%%%%%%%%%%%%
%%%% COMMENCE ICI %%%%
%%%%%%%%%%%%%%%%%%%%%%


%%% EN TÊTE DU CHANT %%%

\beginsong%
% Titre du chant
{L'Hymne Des Femmes}%
% Chanteur(s) ou chanteuse(s)
[by={Écriture collective}]%,
% Date (c'est normal si ça ne s'affiche pas)
%date={1971}]
\MultiwordChords 

%%% PAROLES DU CHANT %%%

% Commentaire sur la chanson
\textnote{Sur l'air du chant des marais}

% Début d'un couplet
\begin{verse} 
No\[Am]us qui sommes \[F]sans \[Am]passé, les femmes
\[Am]Nous qu\[C]i n'av\[E7]ons pas d'\[Am]histoire\[G7]~
\[C]Depuis la n\[C]uit des temps, les femmes
\[Dm]Nous s\[C]ommes le \[E7]conti\[Am]nent noir.
\end{verse}
% Fin d'un couplet

% Début du refrain
\begin{chorus}
\[G7]Levons-\[C]nous femmes escl\[G]aves
\[E7]Et br\[Am]isons nos \[B7]entr\[E7]aves
De\[Am]-bou\[E7]t, de\[Am]-bou\[Am]t, de\[E7]-bou\[Am]t !
\end{chorus}
% Fin du refrain

% Autre couplet, pas besoin de remettre les accords si ce sont les mêmes
\begin{verse}
Asservies, humiliées, les femmes
Achetées, vendues, violées
Dans toutes les maisons, les femmes
Hors du monde reléguées.
\end{verse}

% Pas besoin de remettre le refrain si c'est toujours le même

\begin{verse}
Seules dans notre malheur, les femmes
L'une de l'autre ignorée
Ils nous ont divisées, les femmes
Et de nos soeurs séparées.
\end{verse}

\begin{verse}
Le temps de la colère, les femmes
Notre temps, est arrivé
Connaissons notre force, les femmes
Découvrons-nous des milliers !
\end{verse}

\begin{verse}
Reconnaissons-nous, les femmes
Parlons-nous, regardons-nous,
Ensemble, on nous opprime, les femmes
Ensemble, Révoltons-nous !
\end{verse}

% Cette fois le refrain change alors !
\begin{chorus}
% Bon ça ce n'est pas vraiment dans la chanson, mais pour l'exemple on peut mettre une partie en bis
    \begin{bis}
    Levons-nous femmes esclaves
    Et jouissons sans entraves
    Debout, debout, debout !
    \end{bis}
% Fin de la partie en bis
\end{chorus}



%%%%%%%%%%%
%%% FIN %%%
%%%%%%%%%%%

\endsong
\end{songs}
\end{document}

\end{document}